\documentclass[10pt,
% a4paper,
%twocolumn,
fleqn,
%landscape, 
% papersize,
dvipdfmx,
uplatex
]{jsarticle}



\def\maru#1{\textcircled{\scriptsize#1}}%丸囲み番号

% \RequirePackage[2020/09/30]{platexrelease}

%太字設定
\usepackage[deluxe]{otf}

\usepackage{emathEy}

\usepackage[g]{esvect}

%大きな文字
\usepackage{fix-cm}

%定理環境
\usepackage{emathThm}
%\theoremstyle{boxed}
\theorembodyfont{\normalfont}
\newtheorem{Question}{問題}[subsection]
\newtheorem{Q}{}[subsection]
\newtheorem{question}[Question]{}
\newtheorem{quuestion}{}[subsection]

%セクション，大問番号のデザイン
\renewcommand{\labelenumi}{(\arabic{enumi})}
\renewcommand{\theenumii}{\alph{enumii})}
\renewcommand{\thesection}{第\arabic{section}章}

%用紙サイズの詳細設定
\usepackage{bxpapersize}
\papersizesetup{size={80mm,45mm}}
\usepackage[top=0.7zw,bottom=0truemm,left=3truemm,right=133truemm]{geometry}
\usepackage[dvipdfmx]{graphicx}

%余白など
\usepackage{setspace} % 行間
\setlength{\mathindent}{1zw}
\setlength\parindent{0pt}


%色カラーに関する設定
\usepackage{color}
\definecolor{shiro}{rgb}{0.95703125,0.87109375,0.7421875}
\definecolor{kin}{rgb}{0.95703125,0.87109375,0.7421875}
\definecolor{orange}{rgb}{1,0.7,0.2}
\definecolor{bradorange}{rgb}{1,0.5,0}
\definecolor{pink}{rgb}{0.9176,0.5686,0.5960}
\definecolor{mizu}{rgb}{0.6156,0.8,0.9955}
\color{kin}
% \pagecolor{hukamido}

\usepackage{at}%図の配置
% \usepackage{wallpaper}

\begin{document}



\at(0cm,0cm){\includegraphics[width=8cm,bb=0 0 1920 1080]{./youtube/thumbnails/templates/smart_background/数III微分.jpeg}}
{\color{orange}\bf\boldmath\huge\underline{整数問題への応用}}\vspace{0.3zw}

\huge 
\bf\boldmath 問．$a^{b^2}=b^{a^2}$かつ$a<b$をみたす自然数の組$\left(a,\;b\right)$は存在するか．
\at(7.0cm,0.2cm){\small\color{bradorange}$\overset{\text{数Ⅲ微分}}{\text{典型}}$}


\newpage



\at(0cm,0cm){\includegraphics[width=8cm,bb=0 0 1920 1080]{./youtube/thumbnails/templates/smart_background/数III微分.jpeg}}
{\color{orange}\bf\boldmath\huge\underline{関数の最大最小}}\vspace{0.3zw}

\huge 
\bf\boldmath 問．$-\pi \leqq x\leqq \pi$における$y=2\sin x+\sin 2x$の最大値と最小値を求めよ．
\at(7.0cm,0.2cm){\small\color{bradorange}$\overset{\text{数Ⅲ微分}}{\text{典型}}$}


\newpage



\at(0cm,0cm){\includegraphics[width=8cm,bb=0 0 1920 1080]{./youtube/thumbnails/templates/smart_background/数III微分.jpeg}}
{\color{orange}\bf\boldmath\LARGE\underline{$f''\left(x\right)$を用いた不等式証明}}\vspace{0.3zw}

\LARGE 
\bf\boldmath 問．すべての正の数$x$に対して，$e^x>1+x+\bunsuu{{x^2}}{2}$が成立することを示せ．
\at(7.0cm,0.2cm){\small\color{bradorange}$\overset{\text{数Ⅲ微分}}{\text{典型}}$}


\newpage



\at(0cm,0cm){\includegraphics[width=8cm,bb=0 0 1920 1080]{./youtube/thumbnails/templates/smart_background/数III微分.jpeg}}
{\color{orange}\bf\boldmath\huge\underline{定数分離}}\vspace{0.3zw}

\LARGE 
\bf\boldmath 問．方程式$ax^5-x^2+3=0$が$3$個の異なる実数解をもつような$a$の値の範囲を求めよ．
\at(7.0cm,0.2cm){\small\color{bradorange}$\overset{\text{数Ⅲ微分}}{\text{典型}}$}


\newpage



\at(0cm,0cm){\includegraphics[width=8cm,bb=0 0 1920 1080]{./youtube/thumbnails/templates/smart_background/数III微分.jpeg}}
{\color{orange}\bf\boldmath\huge\underline{凹凸グラフの概形}}\vspace{0.3zw}

\LARGE 
\bf\boldmath 問．関数$f\left(x\right)=\bunsuu{{x^2+2x+2}}{{x+1}}$の増減，極値，グラフの凹凸，漸近線を調べ，グラフの概形をかけ．
\at(7.0cm,0.2cm){\small\color{bradorange}$\overset{\text{数Ⅲ微分}}{\text{典型}}$}


\newpage



\at(0cm,0cm){\includegraphics[width=8cm,bb=0 0 1920 1080]{./youtube/thumbnails/templates/smart_background/数III微分.jpeg}}
{\color{orange}\bf\boldmath\LARGE\underline{極値$\cdot$変曲点をもつ条件}}\vspace{0.3zw}

\normalsize 
\bf\boldmath 問．$f\left(x\right)=\left(x^2+a\right)e^x$とする．ただし，$a$は定数とする．\\
(1)  関数$f\left(x\right)$が極値をもたないような$a$の値の範囲を求めよ．\\
(2)  曲線$y=f\left(x\right)$が変曲点をもつような$a$の値の範囲を求めよ\\

\at(7.0cm,0.2cm){\small\color{bradorange}$\overset{\text{数Ⅲ微分}}{\text{典型}}$}


\newpage



\at(0cm,0cm){\includegraphics[width=8cm,bb=0 0 1920 1080]{./youtube/thumbnails/templates/smart_background/確率.jpeg}}
{\color{orange}\bf\boldmath\huge\underline{確率の最大化}}\vspace{0.3zw}

\Large 
\bf\boldmath 問．$O$さんが各問題に正解する確率は$\bunsuu{{99}}{{100}}$である．$O$さんが$3$問違えるまで問題を解き続けるとき，$n$問目で終わる確率$P_n$が最大となる$n$を求めよ．
\at(7.2cm,0.2cm){\small\color{bradorange}$\overset{\text{確率}}{\text{典型}}$}


\newpage



\at(0cm,0cm){\includegraphics[width=8cm,bb=0 0 1920 1080]{./youtube/thumbnails/templates/smart_background/確率.jpeg}}
{\color{orange}\bf\boldmath\huge\underline{条件付き確率 Lv.3 }}\vspace{0.3zw}

\large 
\bf\boldmath 問．$5$回に$1$回の割合で，帽子を忘れる癖のある$N$君が，正月に$A,\;B,\;C$の$3$軒を順に年始廻りをして家に帰ったとき，帽子を忘れてきたことに気づいた．家$B$に忘れてきた確率を求めよ．
\at(7.2cm,0.2cm){\small\color{bradorange}$\overset{\text{確率}}{\text{典型}}$}


\newpage



\at(0cm,0cm){\includegraphics[width=8cm,bb=0 0 1920 1080]{./youtube/thumbnails/templates/smart_background/確率.jpeg}}
{\color{orange}\bf\boldmath\huge\underline{条件付き確率 Lv.2 }}\vspace{0.3zw}

\normalsize 
\bf\boldmath 問．$2$つの箱$A,\;B$があり，$A$には赤玉$4$個と白玉$1$個，$B$には赤玉$2$個と白玉$3$個が入っている．サイコロを振り，$1$の目が出れば$A,\;$他の目が出れば$B$を選び，選んだ箱から玉を$1$個取り出す．取り出した玉が赤であるとき，箱$A$が選ばれていた確率を求めよ．

\at(7.2cm,0.2cm){\small\color{bradorange}$\overset{\text{確率}}{\text{典型}}$}


\newpage



\at(0cm,0cm){\includegraphics[width=8cm,bb=0 0 1920 1080]{./youtube/thumbnails/templates/smart_background/確率.jpeg}}
{\color{orange}\bf\boldmath\LARGE\underline{条件付き確率 Lv.1 }}\vspace{0.3zw}

\Large 
\bf\boldmath 問．当たり$2$本，ハズレ$3$本入った箱からくじを$1$本取り出し，それを元に戻さずにもう$1$本取り出す．$2$本目が当たりだったとき，$1$本目も当たりである確率を求めよ．
\at(7.2cm,0.2cm){\small\color{bradorange}$\overset{\text{確率}}{\text{典型}}$}


\newpage



\at(0cm,0cm){\includegraphics[width=8cm,bb=0 0 1920 1080]{./youtube/thumbnails/templates/smart_background/確率.jpeg}}
{\color{orange}\bf\boldmath\huge\underline{独立反復試行 Lv.2 }}\vspace{0.3zw}

\normalsize 
\bf\boldmath 問．$A,\;B$の$2$人が繰り返し試合を行う．各試合において，$A$が勝つ確率は$p,\;B$が勝つ確率は$q$で，引き分けはない．先に$3$勝リードした方が優勝とするとき，次の確率を求めよ．\\
(1)  $5$試合目に$A$が優勝を決める確率\\
(2)  $9$試合目に$A$が優勝を決める確率\\

\at(7.2cm,0.2cm){\small\color{bradorange}$\overset{\text{確率}}{\text{典型}}$}


\newpage



\at(0cm,0cm){\includegraphics[width=8cm,bb=0 0 1920 1080]{./youtube/thumbnails/templates/smart_background/確率.jpeg}}
{\color{orange}\bf\boldmath\huge\underline{独立反復試行 Lv.1 }}\vspace{0.3zw}

\normalsize 
\bf\boldmath 問．$A,\;B$の$2$人が繰り返し試合を行う．各試合において，$A$が勝つ確率は$p,\;B$が勝つ確率は$q$で，引き分けはない．先に$4$勝した方が優勝とするとき，次の確率を求めよ．\\
(1)  $6$試合目に$A$が優勝を決める確率\\
(2)  $6$試合目に優勝者が決まる確率\\

\at(7.2cm,0.2cm){\small\color{bradorange}$\overset{\text{確率}}{\text{典型}}$}


\newpage



\at(0cm,0cm){\includegraphics[width=8cm,bb=0 0 1920 1080]{./youtube/thumbnails/templates/smart_background/確率.jpeg}}
{\color{orange}\bf\boldmath\LARGE\underline{ランダムウォーク Lv.2 }}\vspace{0.3zw}

\normalsize 
\bf\boldmath 問．数直線上の動点$\text{P}$を，コインを投げて
表が出れば正の向きに$1$だけ移動させ，
裏が出れば負の向きに$1$だけ移動させる．
原点$\text{O}$から出発して，コインを${10}$回投げた後に点$\text{P}$が初めて原点に戻る確率を求めよ．
\at(7.2cm,0.2cm){\small\color{bradorange}$\overset{\text{確率}}{\text{典型}}$}


\newpage



\at(0cm,0cm){\includegraphics[width=8cm,bb=0 0 1920 1080]{./youtube/thumbnails/templates/smart_background/確率.jpeg}}
{\color{orange}\bf\boldmath\LARGE\underline{ランダムウォーク Lv.1 }}\vspace{0.3zw}

\normalsize 
\bf\boldmath 問．数直線上の動点$\text{P}$を，コインを投げて
表が出れば正の向きに$1$だけ移動さ，
裏が出れば負の向きに$1$だけ移動させる．
原点$\text{O}$から出発して，コインを${10}$回投げた後の点$\text{P}$が正の部分にある確率を求めよ．
\at(7.2cm,0.2cm){\small\color{bradorange}$\overset{\text{確率}}{\text{典型}}$}


\newpage



\at(0cm,0cm){\includegraphics[width=8cm,bb=0 0 1920 1080]{./youtube/thumbnails/templates/smart_background/確率.jpeg}}
{\color{orange}\bf\boldmath\huge\underline{全体像を見る}}\vspace{0.3zw}

\Large 
\bf\boldmath 問．赤玉$3$個，白玉$3$個，青玉$3$個が入っている袋から$3$個の玉を取り出すとき，玉の色が$2$種類になる確率を求めよ．
\at(7.2cm,0.2cm){\small\color{bradorange}$\overset{\text{確率}}{\text{典型}}$}


\newpage



\at(0cm,0cm){\includegraphics[width=8cm,bb=0 0 1920 1080]{./youtube/thumbnails/templates/smart_background/確率.jpeg}}
{\color{orange}\bf\boldmath\huge\underline{余事象の利用}}\vspace{0.3zw}

\large 
\bf\boldmath 問．$1$から$8$までの数の書かれた$8$枚のカードから$3$枚のカードを取り出すとき，次の確率を求めよ．\\
(1)  $3$数の和が${18}$以下となる確率\\
(2)  $3$数の積が$4$の倍数となる確率\\

\at(7.2cm,0.2cm){\small\color{bradorange}$\overset{\text{確率}}{\text{典型}}$}


\newpage



\at(0cm,0cm){\includegraphics[width=8cm,bb=0 0 1920 1080]{./youtube/thumbnails/templates/smart_background/確率.jpeg}}
{\color{orange}\bf\boldmath\large\underline{非復元抽出$〜$引いたくじは戻さない$〜$}}\vspace{0.3zw}

\Large 
\bf\boldmath 問．当たり$3$本，はずれ$7$本のくじから$4$本を引くとき，$2$本だけ当たりくじを引く確率を求めよ．ただし，引いたくじは戻さないとする．
\at(7.2cm,0.2cm){\small\color{bradorange}$\overset{\text{確率}}{\text{典型}}$}


\newpage



\at(0cm,0cm){\includegraphics[width=8cm,bb=0 0 1920 1080]{./youtube/thumbnails/templates/smart_background/確率.jpeg}}
{\color{orange}\bf\boldmath\huge\underline{同基準}}\vspace{0.3zw}

\LARGE 
\bf\boldmath 問．トランプのスペード${13}$枚を一列に並べるとき，絵札がすべて隣り合う確率を求めよ．
\at(7.2cm,0.2cm){\small\color{bradorange}$\overset{\text{確率}}{\text{典型}}$}


\newpage



\at(0cm,0cm){\includegraphics[width=8cm,bb=0 0 1920 1080]{./youtube/thumbnails/templates/smart_background/確率.jpeg}}
{\color{orange}\bf\boldmath\huge\underline{等確率}}\vspace{0.3zw}

\LARGE 
\bf\boldmath 問．区別のない$3$個のサイコロを投げるとき，出た目の和が$5$となる確率を求めよ．
\at(7.2cm,0.2cm){\small\color{bradorange}$\overset{\text{確率}}{\text{典型}}$}


\newpage



\at(0cm,0cm){\includegraphics[width=8cm,bb=0 0 1920 1080]{./youtube/thumbnails/templates/smart_background/数I数と式.jpeg}}
{\color{orange}\bf\boldmath\huge\underline{根号の計算}}\vspace{0.3zw}

\huge 
\bf\boldmath 問．$\alpha =2+\sqrt 3$のとき，$\alpha ^4+\alpha ^3+\alpha ^2+\alpha +1$の値を求めよ．
\at(6.8cm,0.2cm){\small\color{bradorange}$\overset{\text{数Ｉ数と式}}{\text{典型}}$}


\newpage



\at(0cm,0cm){\includegraphics[width=8cm,bb=0 0 1920 1080]{./youtube/thumbnails/templates/smart_background/黒板風.jpeg}}
{\color{orange}\bf\boldmath\huge\underline{$2$重根号}}\vspace{0.3zw}

\HUGE 
\bf\boldmath 問．$\sqrt {2+\sqrt 3}$の$2$重根号をはずせ．
\at(6.6cm,0.2cm){\small\color{bradorange}$\overset{\text{別解研究入門}}{\text{}}$}


\newpage



\at(0cm,0cm){\includegraphics[width=8cm,bb=0 0 1920 1080]{./youtube/thumbnails/templates/smart_background/２次関数.jpeg}}
{\color{orange}\bf\boldmath\huge\underline{解の配置 Lv.3 }}\vspace{0.3zw}

\LARGE 
\bf\boldmath 問．$2$次方程式$x^2-\left(k+4\right)x-\bunsuu{k}{2}+4=0$が$1<x<4$に少なくとも$1$つの実数解をもつような実数$k$の値の範囲を求めよ．
\at(7.0cm,0.2cm){\small\color{bradorange}$\overset{\text{２次関数}}{\text{典型}}$}


\newpage



\at(0cm,0cm){\includegraphics[width=8cm,bb=0 0 1920 1080]{./youtube/thumbnails/templates/smart_background/２次関数.jpeg}}
{\color{orange}\bf\boldmath\huge\underline{解の配置 Lv.2 }}\vspace{0.3zw}

\LARGE 
\bf\boldmath 問．$2$次方程式$ax^2-\left(a-1\right)x-a+1=0$が$-1<x<1$と$3<x<4$にそれぞれ$1$つの実数解を持つような定数$a$の値の範囲を求めよ．
\at(7.0cm,0.2cm){\small\color{bradorange}$\overset{\text{２次関数}}{\text{典型}}$}


\newpage



\at(0cm,0cm){\includegraphics[width=8cm,bb=0 0 1920 1080]{./youtube/thumbnails/templates/smart_background/２次関数.jpeg}}
{\color{orange}\bf\boldmath\huge\underline{解の配置 Lv.1 }}\vspace{0.3zw}

\large 
\bf\boldmath 問．$2$次方程式$x^2+2ax+3-2a=0$が次のような条件を満たすような実数$a$の値の範囲を求めよ．\\
(1)  符号の異なる$2$解をもつ\\
(2)  正の解をもつ\\

\at(7.0cm,0.2cm){\small\color{bradorange}$\overset{\text{２次関数}}{\text{典型}}$}


\newpage



\at(0cm,0cm){\includegraphics[width=8cm,bb=0 0 1920 1080]{./youtube/thumbnails/templates/smart_background/数I数と式.jpeg}}
{\color{orange}\bf\boldmath\LARGE\underline{$2$つの文字含む因数分解}}\vspace{0.3zw}

\LARGE 
\bf\boldmath 問．$2$元$2$次式

\vspace{0.3zw}
\hspace{0.5zw}$6x^2+5xy+y^2-7x-3y+2\vspace{0.3zw}$


を因数分解せよ．
\at(6.8cm,0.2cm){\small\color{bradorange}$\overset{\text{数Ｉ数と式}}{\text{典型}}$}


\newpage



\at(0cm,0cm){\includegraphics[width=8cm,bb=0 0 1920 1080]{./youtube/thumbnails/templates/smart_background/２次関数.jpeg}}
{\color{orange}\bf\boldmath\large\underline{条件式付き$2$変数関数の最大最小 Lv.2 }}\vspace{0.3zw}

\LARGE 
\bf\boldmath 問．$2x^2+3y^2=8$のとき，$4x+3y^2$の最大値および最小値を求めよ．
\at(7.0cm,0.2cm){\small\color{bradorange}$\overset{\text{２次関数}}{\text{典型}}$}


\newpage



\at(0cm,0cm){\includegraphics[width=8cm,bb=0 0 1920 1080]{./youtube/thumbnails/templates/smart_background/２次関数.jpeg}}
{\color{orange}\bf\boldmath\LARGE\underline{複$2$次$4$次関数の最小値}}\vspace{0.3zw}

\huge 
\bf\boldmath 問．関数$y=x^4+6x^2+{10}$の最小値を求めよ．
\at(7.0cm,0.2cm){\small\color{bradorange}$\overset{\text{２次関数}}{\text{典型}}$}


\newpage



\at(0cm,0cm){\includegraphics[width=8cm,bb=0 0 1920 1080]{./youtube/thumbnails/templates/smart_background/２次関数.jpeg}}
{\color{orange}\bf\boldmath\LARGE\underline{独立$2$変数関数の最小値}}\vspace{0.3zw}

\LARGE 
\bf\boldmath 問．$x$と$y$が互いに関係なく変化するとき，$P=x^2+2y^2-2xy+2x+3$の最小値とそのときの$x,\;y$の値を求めよ．
\at(7.0cm,0.2cm){\small\color{bradorange}$\overset{\text{２次関数}}{\text{典型}}$}


\newpage



\at(0cm,0cm){\includegraphics[width=8cm,bb=0 0 1920 1080]{./youtube/thumbnails/templates/smart_background/２次関数.jpeg}}
{\color{orange}\bf\boldmath\Large\underline{$2$次関数の最大値$M$の最小値}}\vspace{0.3zw}

\normalsize 
\bf\boldmath 問．$a$を与えられた定数として$x$の$2$次関数$y=-x^2+4ax+4a$を考え，その最大値を$M$とする．\\
(1)  $M$を$a$の式で表せ．\\
(2)  $M$を最小とする$a$の値を求めよ．また，そのときの$M$の値を求めよ．\\

\at(7.0cm,0.2cm){\small\color{bradorange}$\overset{\text{２次関数}}{\text{典型}}$}


\newpage



\at(0cm,0cm){\includegraphics[width=8cm,bb=0 0 1920 1080]{./youtube/thumbnails/templates/smart_background/２次関数.jpeg}}
{\color{orange}\bf\boldmath\Large\underline{$2$次関数の最大最小から係数決定}}\vspace{0.3zw}

\LARGE 
\bf\boldmath 問．関数$y=ax^2-2ax+b\left(0\leqq x\leqq 3\right)$の最大値が$9$で，最小値が$1$であるとき，定数$a,\;b$の値を求めよ．
\at(7.0cm,0.2cm){\small\color{bradorange}$\overset{\text{２次関数}}{\text{典型}}$}


\newpage



\at(0cm,0cm){\includegraphics[width=8cm,bb=0 0 1920 1080]{./youtube/thumbnails/templates/smart_background/確率.jpeg}}
{\color{orange}\bf\boldmath\Large\underline{サイコロの目の最大値と最小値}}\vspace{0.3zw}

\large 
\bf\boldmath 問．サイコロを$n$回振り，出た目の最大値を$M,\;$最小値を$m$とする．\\
(1)  $M=5$となる確率を求めよ．\\
(2)  $M=5,\;m=2$となる確率を求めよ．\\
(3)  $M-m=3$となる確率を求めよ．\\

\at(7.2cm,0.2cm){\small\color{bradorange}$\overset{\text{確率}}{\text{典型}}$}


\newpage



\at(0cm,0cm){\includegraphics[width=8cm,bb=0 0 1920 1080]{./youtube/thumbnails/templates/smart_background/確率.jpeg}}
{\color{orange}\bf\boldmath\huge\underline{サイコロの目の積}}\vspace{0.3zw}

\normalsize 
\bf\boldmath 問．サイコロを$n$回振り，出た目のすべての積を$X$とするとき，\\
(1)  $X$が偶数である確率を求めよ．\\
(2)  $X$が$6$の倍数である確率を求めよ．\\
(3)  $X$が$4$の倍数である確率を求めよ．\\
(4)  $X$が${12}$の倍数である確率を求めよ．\\

\at(7.2cm,0.2cm){\small\color{bradorange}$\overset{\text{確率}}{\text{典型}}$}


\newpage



\at(0cm,0cm){\includegraphics[width=8cm,bb=0 0 1920 1080]{./youtube/thumbnails/templates/smart_background/数I数と式.jpeg}}
{\color{orange}\bf\boldmath\Large\underline{絶対値を含む方程式$\cdot$不等式}}\vspace{0.3zw}

\large 
\bf\boldmath 問．次の方程式$\cdot$不等式を解け．\\
(1)  $\zettaiti{x+3}=4x$\\
(2)  $\zettaiti{2x-1}\leqq x+3$\\
(3)  $\zettaiti{x}+\zettaiti{x-1}=3x$\\
(4)  $\zettaiti{x}+\zettaiti{x-1}>3x$\\

\at(6.8cm,0.2cm){\small\color{bradorange}$\overset{\text{数Ｉ数と式}}{\text{典型}}$}


\newpage



\at(0cm,0cm){\includegraphics[width=8cm,bb=0 0 1920 1080]{./youtube/thumbnails/templates/smart_background/数I数と式.jpeg}}
{\color{orange}\bf\boldmath\huge\underline{$1$次不等式の整数解}}\vspace{0.3zw}

\small 
\bf\boldmath 問．次の問いに答えよ．\\
(1)  不等式$\bunsuu{x}{2}+4<\bunsuu{{2x+7}}{3}$を満たす最小の整数
$x$を求めよ．\\
(2)  $x$の不等式$2x+a>5\left(x-1\right)$を満たす$x$のうち，最大の整数が$4$であるとき，定数$a$の値の範囲を求めよ．\\
(3)  $x$の連立不等式$7x-5>{13}-2x,\;x+a\geqq 3x+5$を満たす整数$x$がちょうど$5$個存在するとき，定数$a$の値の範囲を求めよ．\\

\at(6.8cm,0.2cm){\small\color{bradorange}$\overset{\text{数Ｉ数と式}}{\text{典型}}$}


\newpage



\at(0cm,0cm){\includegraphics[width=8cm,bb=0 0 1920 1080]{./youtube/thumbnails/templates/smart_background/数I数と式.jpeg}}
{\color{orange}\bf\boldmath\Large\underline{係数に文字を含む$1$次不等式}}\vspace{0.3zw}

\Large 
\bf\boldmath 問．次の不等式を解け．ただし，$a$は定数とする．\\
(1)  $ax=2\left(x+a\right)$\\
(2)  $ax<x+2$\\
(3)  $ax+1>x+a^2$\\

\at(6.8cm,0.2cm){\small\color{bradorange}$\overset{\text{数Ｉ数と式}}{\text{典型}}$}


\newpage



\at(0cm,0cm){\includegraphics[width=8cm,bb=0 0 1920 1080]{./youtube/thumbnails/templates/smart_background/２次関数.jpeg}}
{\color{orange}\bf\boldmath\large\underline{区間が動く$2$次関数の最大値$\cdot$最小値}}\vspace{0.3zw}

\LARGE 
\bf\boldmath 問．$a$は定数とする．関数$y=x^2-2x+2$の
$a\leqq x\leqq a+2$における最大値と最小値を求めよ．
\at(7.0cm,0.2cm){\small\color{bradorange}$\overset{\text{２次関数}}{\text{典型}}$}


\newpage



\at(0cm,0cm){\includegraphics[width=8cm,bb=0 0 1920 1080]{./youtube/thumbnails/templates/smart_background/２次関数.jpeg}}
{\color{orange}\bf\boldmath\Large\underline{軸が動く$2$次関数の最大値$\cdot$最小値}}\vspace{0.3zw}

\LARGE 
\bf\boldmath 問．$a$は定数とする．関数$y=x^2-2ax+3a$の$0\leqq x\leqq 4$における最大値と最小値を求めよ．$a$は定数とする．
\at(7.0cm,0.2cm){\small\color{bradorange}$\overset{\text{２次関数}}{\text{典型}}$}


\newpage



\at(0cm,0cm){\includegraphics[width=8cm,bb=0 0 1920 1080]{./youtube/thumbnails/templates/smart_background/場合の数.jpeg}}
{\color{orange}\bf\boldmath\Large\underline{同じものを含む円順列 Lv.2 }}\vspace{0.3zw}

\huge 
\bf\boldmath 問．赤玉$4$個，白玉$2$個，黒玉$2$個を円形に並べる方法は何通りか．
\at(7.0cm,0.2cm){\small\color{bradorange}$\overset{\text{場合の数}}{\text{典型}}$}


\newpage



\at(0cm,0cm){\includegraphics[width=8cm,bb=0 0 1920 1080]{./youtube/thumbnails/templates/smart_background/場合の数.jpeg}}
{\color{orange}\bf\boldmath\Large\underline{同じものを含む円順列 Lv.1 }}\vspace{0.3zw}

\huge 
\bf\boldmath 問．赤玉$4$個，白玉$3$個，黒玉$1$個を円形に並べる方法は何通りか．
\at(7.0cm,0.2cm){\small\color{bradorange}$\overset{\text{場合の数}}{\text{典型}}$}


\newpage



\at(0cm,0cm){\includegraphics[width=8cm,bb=0 0 1920 1080]{./youtube/thumbnails/templates/smart_background/場合の数.jpeg}}
{\color{orange}\bf\boldmath\huge\underline{重複組合せ Lv.2 }}\vspace{0.3zw}

\LARGE 
\bf\boldmath 問．${10}$個のミカンを$A,\;B,\;C$の$3$人配る．$1$つももらわない人がいても良い場合，配り方は何通りあるか．
\at(7.0cm,0.2cm){\small\color{bradorange}$\overset{\text{場合の数}}{\text{典型}}$}


\newpage



\at(0cm,0cm){\includegraphics[width=8cm,bb=0 0 1920 1080]{./youtube/thumbnails/templates/smart_background/場合の数.jpeg}}
{\color{orange}\bf\boldmath\huge\underline{重複組合せ Lv.1 }}\vspace{0.3zw}

\LARGE 
\bf\boldmath 問．${10}$個のミカンを$A,\;B,\;C$の$3$人配る．どの人も少なくとも$1$個はもらう場合，配り方は何通りあるか．
\at(7.0cm,0.2cm){\small\color{bradorange}$\overset{\text{場合の数}}{\text{典型}}$}


\newpage



\at(0cm,0cm){\includegraphics[width=8cm,bb=0 0 1920 1080]{./youtube/thumbnails/templates/smart_background/場合の数.jpeg}}
{\color{orange}\bf\boldmath\huge\underline{組分け Lv.3 }}\vspace{0.3zw}

\huge 
\bf\boldmath 問．${12}$人を$3$つの$4$人組に分ける方法は何通りか．
\at(7.0cm,0.2cm){\small\color{bradorange}$\overset{\text{場合の数}}{\text{典型}}$}


\newpage



\at(0cm,0cm){\includegraphics[width=8cm,bb=0 0 1920 1080]{./youtube/thumbnails/templates/smart_background/場合の数.jpeg}}
{\color{orange}\bf\boldmath\huge\underline{組分け Lv.2 }}\vspace{0.3zw}

\huge 
\bf\boldmath 問．${12}$人を$2$つの$6$人組に分ける方法は何通りか．
\at(7.0cm,0.2cm){\small\color{bradorange}$\overset{\text{場合の数}}{\text{典型}}$}


\newpage



\at(0cm,0cm){\includegraphics[width=8cm,bb=0 0 1920 1080]{./youtube/thumbnails/templates/smart_background/場合の数.jpeg}}
{\color{orange}\bf\boldmath\huge\underline{組分け Lv.1 }}\vspace{0.3zw}

\huge 
\bf\boldmath 問．${12}$人を$5$人組と$4$人組と$3$人組に分ける方法は何通りか．
\at(7.0cm,0.2cm){\small\color{bradorange}$\overset{\text{場合の数}}{\text{典型}}$}


\newpage



\at(0cm,0cm){\includegraphics[width=8cm,bb=0 0 1920 1080]{./youtube/thumbnails/templates/smart_background/場合の数.jpeg}}
{\color{orange}\bf\boldmath\huge\underline{組分け Lv.0 }}\vspace{0.3zw}

\huge 
\bf\boldmath 問．${12}$人を$8$人組と$4$人組に分ける方法は何通りか．
\at(7.0cm,0.2cm){\small\color{bradorange}$\overset{\text{場合の数}}{\text{典型}}$}


\newpage



\at(0cm,0cm){\includegraphics[width=8cm,bb=0 0 1920 1080]{./youtube/thumbnails/templates/smart_background/場合の数.jpeg}}
{\color{orange}\bf\boldmath\huge\underline{部屋分け Lv.2 }}\vspace{0.3zw}

\LARGE 
\bf\boldmath 問．空室は作らないものとする．$6$人を$A,\;B,\;C$の$3$部屋に分ける方法は何通りあるか．
\at(7.0cm,0.2cm){\small\color{bradorange}$\overset{\text{場合の数}}{\text{典型}}$}


\newpage



\at(0cm,0cm){\includegraphics[width=8cm,bb=0 0 1920 1080]{./youtube/thumbnails/templates/smart_background/場合の数.jpeg}}
{\color{orange}\bf\boldmath\huge\underline{部屋分け Lv.1 }}\vspace{0.3zw}

\LARGE 
\bf\boldmath 問．空室は作らないものとする．$6$人を$A,\;B$の$2$部屋に分ける方法は何通りあるか．
\at(7.0cm,0.2cm){\small\color{bradorange}$\overset{\text{場合の数}}{\text{典型}}$}


\newpage



\at(0cm,0cm){\includegraphics[width=8cm,bb=0 0 1920 1080]{./youtube/thumbnails/templates/smart_background/集合と命題.jpeg}}
{\color{orange}\bf\boldmath\huge\underline{背理法証明 Lv.2 }}\vspace{0.3zw}

\large 
\bf\boldmath 問．有理数$a,\;b$のうち少なくとも$1$つが$0$でないならば，$a\sqrt 2+b\sqrt 3$は無理数であることを証明せよ．ただし，平方数でない正の整数$m$に対して，$\sqrt m$が無理数であることを前提としてよい．
\at(6.8cm,0.2cm){\small\color{bradorange}$\overset{\text{集合と命題}}{\text{典型}}$}


\newpage



\at(0cm,0cm){\includegraphics[width=8cm,bb=0 0 1920 1080]{./youtube/thumbnails/templates/smart_background/集合と命題.jpeg}}
{\color{orange}\bf\boldmath\huge\underline{背理法証明 Lv.1 }}\vspace{0.3zw}

\Large 
\bf\boldmath 問．$\sqrt 2+\sqrt 3$は無理数であることを証明せよ．ただし，平方数でない正の整数$m$に対して，$\sqrt m$が無理数であることを前提としてよい．
\at(6.8cm,0.2cm){\small\color{bradorange}$\overset{\text{集合と命題}}{\text{典型}}$}


\newpage



\at(0cm,0cm){\includegraphics[width=8cm,bb=0 0 1920 1080]{./youtube/thumbnails/templates/smart_background/集合と命題.jpeg}}
{\color{orange}\bf\boldmath\Large\underline{$\sqrt 6$は無理数であることの証明}}\vspace{0.3zw}

\huge 
\bf\boldmath 問．$\sqrt 2$は無理数であることを証明せよ．
\at(6.8cm,0.2cm){\small\color{bradorange}$\overset{\text{集合と命題}}{\text{典型}}$}


\newpage



\at(0cm,0cm){\includegraphics[width=8cm,bb=0 0 1920 1080]{./youtube/thumbnails/templates/smart_background/集合と命題.jpeg}}
{\color{orange}\bf\boldmath\Large\underline{$\sqrt 2$は無理数であることの証明}}\vspace{0.3zw}

\huge 
\bf\boldmath 問．$\sqrt 6$は無理数である．ことを証明せよ．
\at(6.8cm,0.2cm){\small\color{bradorange}$\overset{\text{集合と命題}}{\text{典型}}$}


\newpage



\at(0cm,0cm){\includegraphics[width=8cm,bb=0 0 1920 1080]{./youtube/thumbnails/templates/smart_background/集合と命題.jpeg}}
{\color{orange}\bf\boldmath\huge\underline{対偶証明法 Lv.3 }}\vspace{0.3zw}

\LARGE 
\bf\boldmath 問．整数$n$について，$n^2$が$6$の倍数ならば，$n$も$6$の倍数であることを証明せよ．
\at(6.8cm,0.2cm){\small\color{bradorange}$\overset{\text{集合と命題}}{\text{典型}}$}


\newpage



\at(0cm,0cm){\includegraphics[width=8cm,bb=0 0 1920 1080]{./youtube/thumbnails/templates/smart_background/集合と命題.jpeg}}
{\color{orange}\bf\boldmath\huge\underline{対偶証明法 Lv.2 }}\vspace{0.3zw}

\LARGE 
\bf\boldmath 問．整数$n$について，$n^2$が$3$の倍数ならば，$n$も$3$の倍数であることを証明せよ．
\at(6.8cm,0.2cm){\small\color{bradorange}$\overset{\text{集合と命題}}{\text{典型}}$}


\newpage



\at(0cm,0cm){\includegraphics[width=8cm,bb=0 0 1920 1080]{./youtube/thumbnails/templates/smart_background/集合と命題.jpeg}}
{\color{orange}\bf\boldmath\huge\underline{対偶証明法 Lv.1 }}\vspace{0.3zw}

\LARGE 
\bf\boldmath 問．整数$n$について，$n^2$が$2$の倍数ならば，$n$も$2$の倍数であることを証明せよ．
\at(6.8cm,0.2cm){\small\color{bradorange}$\overset{\text{集合と命題}}{\text{典型}}$}


\end{document}

